Sam wanted a proper table and there was not one to be had that did not already have four other purposes, so in the end they used the map table in the small hall behind the kitchens, with the burial council’s chalk marks for the Dragonpit still faintly visible under a candle where nobody had thought to clean it, and Tyrion Lannister thought, sitting down at it, that there was something almost decent about hearing the end of one secret in the same room where they had settled how to burn thirty thousand bodies. Both things, in the end, were about what a government could be trusted to do with the dead.

Sam had three pens laid out, which struck Tyrion as excessive until he remembered that Sam Tarly did not trust anything with only one of a thing that mattered.

“I want it read back to me before it goes anywhere,” Howland Reed said, before anyone had asked him to begin. He had washed, and been given a dry coat, and looked no less tired for either. “Every part of it. If I’ve said a thing wrong I’ll want the chance to say it right before it’s a fact for a hundred years.”

“That’s the arrangement,” Sam said. “I write it as you speak it, and I read it back at the end, and you tell me where I’ve got it wrong, and I cross out nothing — I write the correction beside it, in the margin, with the date, so both stand and anyone after us can see which came first.”

Howland Reed looked at him for a moment with something that was not quite approval and was close to it.

“That’ll do,” he said.

\sceneBreak

Sansa Stark sat at the end of the table with her hands folded and did not say why she had asked to be there, and nobody asked her to leave, and Tyrion understood without being told that this was not a courtesy anyone in the room was prepared to argue about.

“You knew my father before he was Lord of anything,” she said, when Howland had settled.

“I knew him when he was a boy who thought the Vale was the whole of the world, aye. Before Robert’s war, before any of it.” Howland’s hands were flat on the table, the way a man’s hands go when he has decided a thing is going to be difficult and has stopped pretending otherwise to himself. “I’ll not tell it prettier than it was, my lady, and I’d ask you to hear that plainly before I start, because it wasn’t pretty, and your father carried the ugliest part of it alone for years so that nobody else would have to.”

“I’ve heard the shape of it already,” Sansa said. “The annulment. The marriage. The boy.”

“You’ve heard the shape,” Howland agreed. “I was there for the part with no shape in it at all. I’d like to give you that, if you’ll have it, because I think it’s owed, and because I don’t imagine I’ll pass this way again to give it a second time.”

\sceneBreak

He told it plainly, the way he had promised, and Tyrion — who had spent a professional lifetime listening to men perform grief for effect and had got very good at hearing the seams in it — heard none anywhere in the telling.

There had been six of them who rode to that tower with Eddard Stark, and three who came away from it walking, and Howland Reed was one of the three, and he did not dwell on the fight itself beyond what the record needed: three men in white cloaks who would not yield the tower though the war they were guarding it for had already been lost three moons past at a river the losing side no longer controlled, and six who came to take it from them, and a morning that ended with three of the six standing and all three of the white cloaks dead, having chosen it over surrender in a manner Howland still, all these years later, did not know whether to call honor or a particular kind of madness that wore honor’s colors.

“Ned went up the tower steps still bleeding from the shoulder and I went after him because I could stand and he couldn’t be stopped from going, and there was a woman in the room at the top who I’d never in my life seen and knew inside a heartbeat exactly who she had to be, because there was no other woman in the whole of Westeros that morning that three men would have died in a doorway to keep a stranger from reaching.”

He stopped there, for the length of perhaps four breaths, and nobody in the room moved to fill it.

“She was dying and she knew it and there wasn’t a maester within two days’ ride who could have changed that, and your father knew it too, my lady, the moment he saw her, though I don’t believe he let himself say the word even in his own head until she made him.” Howland’s voice did not change register through any of this, which Tyrion found, of everything in the room, the hardest thing to sit still through. “She had the boy already. Small, and quiet, and I remember thinking he oughtn’t to have been so quiet, given everything that had just happened to bring him into the world, and I’ve wondered since whether he knew, even that young, that quiet was going to be the thing that kept him alive.”

\sceneBreak

“She said more than ‘promise me,’” Howland said. “I’ve heard it told since, in the North, in a shorter shape than it happened, and I understand why — a dying woman’s words get worn smooth by every mouth that repeats them, the way a river stone gets worn smooth, until what’s left fits in one sentence and can be carried easily. But I was in that room and I heard the whole of it, and I’d like it written the whole way, if Master Tarly’s hand will bear it.”

“It’ll bear it,” Sam said.

“She was afraid of Robert. That’s the part everyone tells and it’s true and it’s the smaller part.” Howland’s jaw worked once. “She was more afraid of the other thing. Her prince — Rhaegar — had talked to her, in the months before, about signs. About a promised thing, and a dragon with three heads, and a destiny he’d fixed his whole life to before he ever fixed his eyes on her, and I think — I only think, my lady, I’ll not swear to what was in a dying woman’s heart as though I could read it — I think she’d come, by the end, to be more frightened of what her prince’s dreaming might make of her son than she ever was of a usurper’s soldiers. She didn’t want him raised toward anything. Not a throne, not a prophecy, not a destiny some dead man had already decided for him before he’d so much as opened his eyes. She wanted him raised toward nothing at all, so that whatever he grew into, he’d have grown into it himself, and no one — not Robert, not the dragon dreams, not any of it — would ever be able to say they’d made him.”

The room was very quiet.

“That’s why she made your father swear the whole of it, my lady, and not the half that’s remembered. Not only \textit{keep him safe}. \textit{Keep him ordinary.} Let him be a boy who wondered why his father looked at him strangely sometimes and never troubled to find out why, rather than a prince who knew from the cradle what he was for. Ned Stark gave her both promises in the same breath, and I don’t believe, in all the years after, he ever worked out how to keep the one without also keeping the other, because they were never two promises. They were one, said twice for emphasis, by a woman who didn’t trust that a single telling would hold against the rest of the world’s telling.”

\sceneBreak

Tyrion sat with that for a while and found that the political part of his mind — the part that had been running figures under the whole of the testimony without asking his permission, the part that had already worked out, three sentences back, that this changed the character of the claim from \textit{documented} to \textit{witnessed and sworn}, which was a different order of thing entirely in any court that had ever existed — had gone, for the length of perhaps half a minute, entirely silent.

It was Sansa who spoke into the silence, and what she said was not what Tyrion had expected.

“He used to look at Jon oddly,” she said. “When we were children. Not unkindly — I want that said, because I’ve heard people say it unkindly and it was never that. He’d watch him sometimes across a hall, for a moment too long, and then catch himself and stop, and I never once understood it and I don’t think Jon did either, and I spent a good part of my childhood being cruel to him without knowing that whatever I saw on my father’s face when he watched him wasn’t coldness at all.” She was not crying, which Tyrion thought afterward took a particular kind of strength to manage, sitting where she was sitting. “It was the opposite of what I thought it was. That’s the whole of what you’ve just told me, Lord Reed, and I’d have liked to know it years ago, and I understand entirely why I couldn’t have.”

“Aye,” Howland said. “So would he have liked you to know it. That’s the cruelty in the middle of the whole arrangement, my lady, and I’ll not dress it up as anything else. A promise kept perfectly is a wound that never gets the chance to close, because the whole use of it is that nobody’s allowed to touch it.”

\sceneBreak

“I have to ask you something,” Tyrion said, when the room had sat with that for as long as it needed to, “and I’d rather ask it plainly than have it sit unasked.”

“Ask it.”

“Who else have you told.”

He heard it land the way it had landed when he said the same thing to Garth in this same wing some weeks before, except that this time he was the one who had said it, and he heard, in the half-second after it left him, exactly what he had reached for first, out of everything a man might reach for on hearing a story like the one he had just been given — not grief, not wonder, not even pity for a boy he had known for two years and fought beside twice. Risk. He had reached for risk, the way another man’s hand goes to a blade, because it was the tool that had kept him alive in three courts and it did not know how to stay in its sheath merely because the moment called for something gentler.

Howland Reed looked at him for a long moment, and Tyrion understood, from the look, that he had been caught doing it, and that the crannogman was not going to pretend otherwise out of politeness.

“No one,” Howland said. “Not Ned, in the end — I told him what I’d say to him, once, years back, when he first started wondering aloud whether the boy oughtn’t to be told, and he asked me not to, and I honored that the same as I’ve honored the rest of it. Not my own wife, who I trusted with everything else I owned in this world. Not my daughter, who I’d trust with my life and haven’t, because it was never mine to give away.” He held Tyrion’s eye. “You’re among the first four people who’ve ever heard the whole of it in one telling, my lord, and I notice that the first question out of your mouth wasn’t whether it was true, or what it cost me to carry it, but who else might already have it. I don’t hold that against you. I expect it’s the reason you’re the man sitting in that chair and not somebody softer who’d have been eaten alive in it by now. But I’d ask you to notice it yourself, because I don’t think you have.”

Tyrion did not have an answer for that, and did not try to invent one.

\sceneBreak

He asked, some while later, whether the record might be kept close — not hidden, he was careful to say, only \textit{close}, seen by the council that would sit and by no one else until that council decided what to do with it — and it was Howland who told him no, in the same flat unhurried voice he had used for everything else in the room.

“I didn’t walk two moons to have it locked in a box that opens at your convenience, my lord, however well you mean by the asking.” He said it without heat. “I’ve spent thirty years being the only man alive who held this. I know exactly what a secret does to the man carrying it, because I’ve been that man, and I’ll not hand it to a room of five people and call that an improvement if all I’ve done is make the room bigger. Master Tarly writes it plain, and it goes in whatever this city is building to keep its record, and any man with a legitimate reason to read it may read it, the same as any other page in that book. That was the whole of my condition for speaking at all, and I’ll not soften it now that I’ve spoken.”

“That’s not how documents of this weight are usually handled.”

“No,” Howland agreed. “I don’t imagine it is. I’d say that’s rather the point of handling it differently.”

Sam, who had said very little through the whole of it beyond the scratch of his pen, looked up at that, and Tyrion caught, for just a moment, something in the boy’s face that he had not expected to see and could not immediately name, and named eventually, walking back to his own rooms that night, as recognition — the particular look of a man watching somebody else say, in a few words, a thing he had spent his whole adult life failing to say properly himself.

\sceneBreak

He read it back to Howland twice before the crannogman was satisfied, and corrected three words, all of them small, all of them the kind of correction that meant nothing to anyone in the room and would mean a very great deal to whoever came looking for the exact truth of it in a hundred years — \textit{a doorway} struck for \textit{the doorway}, because there had been more than one and precision mattered; \textit{I think} inserted before the matter of what had been in a dying woman’s heart, because Howland would not have his own guess recorded as her own certain word; and a number corrected where Sam had misheard him, small and precise, the kind of thing nobody in the room but Howland would ever have noticed was wrong.

Then Sam signed it, and dated it, and asked Howland to make his mark beside the date, and Howland did, in a hand Tyrion had not expected — small, careful, and entirely legible, the hand of a man who had learned his letters properly and had simply never had much cause to use them for anyone’s benefit but his own.

\sceneBreak

Tyrion went back to his own room that night and got the leather case out from under the loose board where he had taken, in the last while, to keeping it, and sat with it a long time without opening it.

He had told himself, the night he first put the letters in a case rather than leaving them loose on a table, that the case was for tidiness. He understood now, sitting with it on his knees in a room with the shutter rattling and no smoke left anywhere outside to smell, that this had never been true, and that the true reason had been considerably simpler and less flattering: a case could be locked, and a locked thing felt, for as long as the lock held, like a thing that had not quite happened yet.

Eleven unfinished drafts in a dead man’s hand, saying what Howland Reed had just said aloud in front of three witnesses and a boy with a very good memory for exact wording.

He did not burn them. He had not burned them the first night either, and he had stopped, some while back, pretending to himself that the not-burning was an oversight rather than a decision he simply had not been willing to look at directly.

What he did instead, for the first time since the battle, was take a fresh sheet, and write, in his own hand, a plain account of when he had found the letters, what they said, and why he had kept them from the record for as long as he had — which was, he wrote, because he had been afraid, and because a piece of information that only one man possesses is a piece of leverage that only one man possesses, and he had told himself for some while that he was holding it in trust rather than holding it for himself, and had only that evening, in a room behind the kitchens, been shown plainly what the difference actually looked like when a better man did it properly.

He signed that, too, and dated it, and put both — Varys’s eleven letters and his own single page — into the same case, and did not lock it.

Then he went and found Sam, who had not yet gone up his two hundred and eleven steps for the night, and gave it to him without saying very much at all, and Sam took it the way he took everything now, without drama, and asked only whether Tyrion wanted a receipt.

“No,” Tyrion said. “I don’t imagine I’ll be asking for it back.”

He went to bed that night without the wine, not from any particular resolve, but because he found, getting into it, that he did not especially want it, which had not happened to him in some while, and which he did not examine too closely for fear of discovering it would not happen again.